\normalsize{ \indent
Primero, se investigó sobre los antecedentes del colegio;
para lo cual, fue necesario realizar una primera entrevista
para averiguar detalles históricos del manejo de información.
Teniendo en cuenta que el problema inicial resultaba ser
muy grande para el tiempo que se dispone, se recolectó
más datos de los necesarios. Toda la información respecto a
este asunto se encuentra en la descripción del escenario.
}
\newline
\normalsize{ \indent
El paso siguiente fue establecer los objetivos de las
entrevistas (se han llevado a cabo dos), las cuales tienen
enfoques distintos por la información que se manejaba del
colegio y por el asentamiento del alcance de la solución.
Para la entrevista número uno, se propuso como objetivos:
}
\begin{itemize}
  \item Conocer los antecedentes del colegio con el
  gestión de docentes y alumnos (si bien, esto puede
  resultar recursivo en el trabajo práctico; es la realidad
  de la situación, ya que lo único que se conocía antes de
  la entrevista era el deseo de tener un gestor de legajos
  y formularios).
  \item Informarse sobre los problemas que posee la escuela
  secundaria respecto a legajos y formularios. Además,
  constatar si existe algún tipo de solución sugerida por
  los futuros clientes del sistema.
  \item Comprobar si las sugerencias generadas por el
  entrevistador para el sistema son viables para los
  entrevistados. Es decir, si las ideas de la posible
  solución son convenientes para la solución que se genere
  (al menos, según los clientes del futuro software).
  \item Recolectar datos sobre los procedimientos para
  legajos y formularios.
\end{itemize}
\ \newline
\normalsize{ \indent
Para la entrevista más reciente, con un alcance ya definido;
han surgido nuevos objetivos de la entrevista, los cuales son:
}
\begin{itemize}
  \item Investigar la clasificación de los docentes según el
  horario.
  \item Identificar los derechos que los usuarios finales
  tienen sobre el software.
  \item Recolectar datos para analizar las condiciones
  necesarias para registrar a los docentes.
  \item Establecer las reglas a cumplir sobre las horas
  cátedra, el cambio de horarios y el \gls{reldoc} (tanto
  por licencia como renuncia a las horas).
  \item Averiguar el método de asistencia y calificaciones
  que se encuentra actualmente en la institución.
  \item Definir el carácter de la solución frente a las
  distintas situaciones que se puedan presentar.
\end{itemize}
\ \newline
\normalsize{ \indent
A continuación, debía decidirse a las personas que serían
entrevistadas; sin embargo, es necesario dar contexto del
tiempo en que se iban a llevar a cabo los encuentros. A las
fechas de realización del trabajo práctico integrador de
ingeniería I (del cual se extrae parte de este trabajo), el
mundo no se encontraba en una situación normal; el virus
COVID-19 estaba en circulación, por lo que el personal que
trabajaba era limitado. Teniendo en cuenta esto, por descarte
ha resultado que sólo dos personas estaban disponibles para
realizar la entrevista; además, coincide que ambos sujetos
acceden a la entrevista en simultáneo. Cabe destacar que ellos
pertenecen a Coordinación Escolar, tal y como se pretendía
inicialmente (sin tener en cuenta el asunto de coronavirus);
la razón es que son el sector que administra los asuntos
relacionados a legajos, formularios, información personal,
horarios, relevamientos, entre otros.
}
\newline
\normalsize{ \indent
Luego, el autor de este documento se encargó de preparar a los
entrevistados; dando una guía del material que se podía
solicitar en conjunto con los asuntos de las preguntas a
darse. Entre los documentos recibidos, se encuentran los
requisitos para legajos de docentes y alumnos (impreso),
la declaración jurada de los docentes (tanto en papel como
un archivo digital), algunos formularios (los cuales no son
de interés, a causa del alcance que se tiene en esta primera
etapa) y el archivo Excel de horarios (no era necesario
obtener el de legajos, bastaba con la observación; dado que
es muy poco efectivo y excesivamente complejo de visualizar,
incluso de forma personal y con detenimiento).
}
\newline
\normalsize{ \indent
Por último, fue momento de tomar una postura ante el tipo
de preguntas y su estructura. Sobre el primer tema mencionado
en este párrafo, se distingue el uso de preguntas abiertas,
cerradas y sondeos para ambas entrevistas. Por otra parte,
la estructura elegida para cada entrevista fue distinta; en
gran parte por la información en el momento que se realizó
cada una.
}
\newline
\normalsize{ \indent
En la primera, se optó por una forma de diamante; comenzando
a indagar sobre diseños que puedan ser parte de la solución
(sea tabla, lista u otro), los métodos de comunicación entre
los usuarios (si los hubiera); el propósito de la acción es
dar una motivación a los entrevistados, ya que así parece ser
que la solución no está muy lejana a la fecha del encuentro y
las personas que contestan las preguntas hacen tal acción con
mayor motivación. Después, se hacen las preguntas más abiertas
para obtener un contexto amplio de la situación; averiguando
los procedimientos y las demandas de los clientes del futuro
software. Al final, se refinan detalles sobre lo discutido
anteriormente para completar la información obtenida.
}
\newline
\normalsize{ \indent
En la segunda entrevista, se optó por una estructura de
pirámide; si bien se informó a los entrevistados sobre el
alcance que se optó en la primera etapa, resultaría más
adecuado adentrar al personal de Coordinación en el tema sin
preguntas amplias. En adición, los datos juntados en la primera
entrevista daban respuesta a muchas de las preguntas específicas
de horarios y relevamientos. Entonces, se ubicaron al principio
preguntas específicas sobre los derechos de acciones de los
usuarios, la clasificación de profesores según horarios y las
reglas correspondientes a cada división (que distingue a cada
categoría). Luego, se prosiguió con la investigación sobre las
asistencias y calificaciones de los docentes y alumnos. Para
concluir, se planifica preguntar sobre el procedimiento de
cambio de horario y relevamiento; en conjunto con sus reglas,
las cuales se suponen que son amplias.
}
\newline
\normalsize{ \indent
Para terminar, cabe mencionar que el uso de sondeos resulta
difícil ubicar en momentos específicos de la entrevista; este
tipo de preguntas se realiza respecto a cierto tipo de respuestas
a una cuestión anterior, por lo que el entrevistador debe prestar
atención al utilizar los mismos en tiempo real. Igualmente, se
tiene una noción de las situaciones clave gracias a la
experiencia.
}
